
\section*{1. Local properties of an \(A\)-group locally of finite type}
\phantomsection\label{sec:sga3-VIA-1}
\addcontentsline{toc}{section}{1. Local properties of an
  \texorpdfstring{\(A\)}{A}-group locally of finite type}

We shall first see that, if \(G\) is locally of finite type and flat over
\(A\), one can ``make every closed point of \(G\) strictly rational'' by
means of a finite flat extension of the base.

\medskip
\phantomsection\label{par:sga3-VIA-1.1}%
\noindent\textbf{1.1.}
Unless expressly stated otherwise, from now on we assume that \(G\) is an
\(A\)-group locally of finite type. When \(A\) is a field \(k\), we shall
obtain in Expos\'e~VII\textsubscript{B} very precise results on the local
rings of \(G\).%
\footnote{\label{note:sga3-VIA-10}Editor's note: for example, they are
always complete intersections; see Expos\'e~VII\textsubscript{B}, \SGARefLink{sga3:VIIA:5.5.1}{5.5.1}.
Moreover, if \(\operatorname{char}(k)=0\), then \(G\) is smooth
(Expos\'e~VI\textsubscript{B}, \SGARefLink{cor:sga3-VIB-1.6.1}{1.6.1}; see also
Expos\'e~VII\textsubscript{B}, \SGARefLink{sga3:VIIA:3.3.1}{3.3.1}).}
Here we content ourselves with a few elementary results.

\medskip
\phantomsection\label{prop:sga3-VIA-1.1.1}%
\noindent\textbf{Proposition 1.1.1.}
\textit{Let \(x\) be a point of an \(A\)-group \(G\) locally of finite type
and flat over \(A\). Then the local ring \(\mathcal O_{G,x}\) is
Cohen--Macaulay, and there is a system of parameters
\(a_1,\ldots,a_n\) of \(\mathcal O_{G,x}\) such that
\(\mathcal O_{G,x}/(a_1,\ldots,a_n)\) is a finite flat \(A\)-module
(hence finite free).}

\medskip
We first suppose that \(A\) is equal to its residue field \(k\). It then
suffices to prove that \(\mathcal O_{G,x}\) is Cohen--Macaulay, and we may
restrict to the case in which \(x\) is a closed point
(cf.\ EGA~\(0_{\mathrm{IV}}\), 16.5.13). By
\hyperref[lem:sga3-VIA-1.1.2]{Lemma~1.1.2} below, \(G\) contains a closed
point \(y\) such that \(\mathcal O_{G,y}\) is Cohen--Macaulay. By
SGA~1, I, \S~9, this is equivalent to saying that, for every finite
extension \(K\) of the base field \(k\) and every point \(\overline y\) of
\(\overline G=G\otimes_k K\) lying above \(y\), the local ring
\(\mathcal O_{\overline G,\overline y}\) is Cohen--Macaulay. If \(K\) has
been chosen sufficiently large, that is, if \(K\) contains a normal
extension of \(k\) containing the residue fields \(\kappa(y)\) and
\(\kappa(x)\), then \(\overline y\) is (strictly) rational over \(K\), as
is every point \(\overline x\) of \(\overline G\) lying above \(x\).%
\footnote{\label{note:sga3-VIA-11}Editor's note: indeed, the hypothesis on
\(K\) implies that, for every extension \(L\) of \(K\), every
\(k\)-morphism \(\kappa(x)\to L\) (respectively,
\(\kappa(y)\to L\)) factors through \(K\). Consequently, all the points of
\(G\otimes_k K\) lying above \(x\) or \(y\) have \(K\) as residue field;
that is, they are (strictly) rational over \(K\).}
Since the automorphism
\[
  r_{\overline x}\circ(r_{\overline y})^{-1}
\]
sends \(\overline y\) to \(\overline x\), it follows that
\(\mathcal O_{\overline G,\overline x}\), and hence
\(\mathcal O_{G,x}\) (SGA~1, I, \S~9), are Cohen--Macaulay.

Now let \(A\) again be arbitrary. The preceding argument applies to
\(k\otimes_A G\), so that \(k\otimes_A\mathcal O_{G,x}\) is
Cohen--Macaulay. If \(a_1,\ldots,a_n\) is a sequence of elements of
\(\mathcal O_{G,x}\) whose image in
\(k\otimes_A\mathcal O_{G,x}\) is a system of parameters, it follows from
SGA~1, IV, 5.7, or EGA~\(0_{\mathrm{IV}}\), 15.1.16, that
\(a_1,\ldots,a_n\) is an \(\mathcal O_{G,x}\)-regular sequence and that
\(\mathcal O_{G,x}/(a_1,\ldots,a_n)\) is finite and flat (hence finite and
free) over \(A\).

\medskip
\phantomsection\label{lem:sga3-VIA-1.1.2}%
\noindent\textbf{Lemma 1.1.2.}
\textit{Every nonempty scheme \(X\), locally of finite type over an
Artinian ring \(A\), contains a closed point \(x\) whose local ring is
Cohen--Macaulay.}

\medskip
We may plainly suppose that \(X\) is affine, with algebra \(B\), and argue
by induction on \(\dim X\) (the assertion is clear if \(X\) is discrete,
since all its local rings are then Artinian). Since \(B\) is of finite type
over \(A\), if \(\dim B>0\), then \(B\) contains an element \(a\) that is
not invertible and is not a zero divisor.%
\footnote{\label{note:sga3-VIA-12}Editor's note: indeed, \(B\) is a
Noetherian Jacobson ring (cf.\ \textup{[BAC]}, V, \S~3.4). If every
noninvertible element is a zero divisor, then every prime ideal is an
associated prime ideal of \(B\); in particular, \(B\) has only finitely
many maximal ideals \(\mathfrak m_1,\ldots,\mathfrak m_n\). Since \(B\)
is a Jacobson ring, the intersection of the \(\mathfrak m_i\) is the
nilradical of \(B\). It follows that each \(\mathfrak m_i\) is a minimal
prime ideal of \(B\), and hence that \(\dim B=0\).}
The closed subscheme
\[
  X'=\operatorname{Spec} B/(a)
\]
of \(X\) then has dimension strictly less than \(\dim X\), and by the
induction hypothesis it contains a closed point \(x\) such that
\(\mathcal O_{X',x}\) is Cohen--Macaulay. Since
\(\mathcal O_{X',x}=\mathcal O_{X,x}/(a)\), and \(a\) is neither invertible
nor a zero divisor in \(\mathcal O_{X,x}\), the ring
\(\mathcal O_{X,x}\) is Cohen--Macaulay (see also EGA~IV\textsubscript{2},
6.11.3).

\medskip
\phantomsection\label{prop:sga3-VIA-1.2}%
\noindent\textbf{Proposition 1.2.}
\textit{Let \(A\) be an Artinian local ring, let \(G\) be an \(A\)-group
locally of finite type and flat over \(A\), and let \(x\) be a closed point
of \(G\). Then there exists a local \(A\)-algebra \(A'\), finite and free
over \(A\), such that every point of \(G\otimes_A A'\) lying above \(x\)
is strictly rational over \(A'\).}

\medskip
\noindent\footnote{\label{note:sga3-VIA-13}Editor's note: the original has
been expanded in what follows.}
Indeed, let \(k_1\) be a finite normal extension of \(k\) containing the
residue field \(\kappa(x)\) of \(x\). By \SGARefLink{lem:sga3-V-4.1.2}{Lemma~V.4.1.2}, there exists a
local \(A\)-algebra \(A_1\), finite and free over \(A\), with residue field
\(k_1\). In this case (cf.\ Editor's note~\ref{note:sga3-VIA-11}), all the
points \(g_1,\ldots,g_n\) of \(G\otimes_A A_1\) lying above \(x\in G\)
have \(k_1\) as residue field; that is, \(g_1,\ldots,g_n\) are rational
over \(A_1\) in the sense of Expos\'e~V, \S~4.e).

Let \(B_1,\ldots,B_n\) be the local rings of \(g_1,\ldots,g_n\).
By \hyperref[prop:sga3-VIA-1.1.1]{1.1.1}, the rings
\(B_1,\ldots,B_n\) have quotients \(B'_1,\ldots,B'_n\) that are Artinian
and finite and free over \(A_1\). Set
\[
  A'=B'_1\otimes_{A_1}\cdots\otimes_{A_1}B'_n.
\]
Then \(A'\) is local, finite, and free over \(A_1\), and, for each
\(i=1,\ldots,n\), there are surjective homomorphisms
\phantomsection\label{eq:sga3-VIA-1.2-surjections}%
\[
  B_i\otimes_{A_1}A'
    \twoheadrightarrow B'_i\otimes_{A_1}A'
    \twoheadrightarrow A',
\]
the second induced by the multiplication map
\[
  B'_i\otimes_{A_1}B'_i\twoheadrightarrow B'_i.
\]
Consequently, \(A'\) has the required property.

\medskip
\phantomsection\label{par:sga3-VIA-1.3}%
\noindent\textbf{1.3.}
Let \(e\) be the unit element (or origin) of \(G\), that is, the image of
the unique point of \(\operatorname{Spec}A\) under the unit section
\(\operatorname{Spec}A\to G\). By definition, \(e\) is strictly rational
over \(A\).

\par\medskip
\noindent\phantomsection\label{prop:sga3-VIA-1.3.1}%
\textbf{Proposition 1.3.1.}%
\footnote{\label{note:sga3-VIA-14}Editor's note: \hyperref[concl:sga3-VIA-1.3.1-1]{part~(1)}, which is useful
in \hyperref[ex:sga3-VIA-1.3.2]{Examples~1.3.2} below, has been added.}
\textit{Let \(G\) be a group locally of finite type and flat over an
Artinian ring \(A\), and let \(K\) (respectively, \(\overline k\)) be the
perfect closure (respectively, an algebraic closure) of the residue field
\(k\) of \(A\).}

\begin{enumerate}
\item[(1)]\phantomsection\label{concl:sga3-VIA-1.3.1-1} \textit{For every closed point \(x\) of
\(\overline G=G\otimes_A\overline k\),\footnote{Source note: before the explicit reduction \(A=k\), the re-edition
prints \(k\) as the tensor base here and again in the proof. Both objects
are defined over the Artinian ring \(A\); the later post-reduction tensor
bases remain \(k\).} the local rings
\(\mathcal O_{\overline G,e}\) and \(\mathcal O_{\overline G,x}\) are
isomorphic. In particular, the tangent spaces \(T_e\overline G\) and
\(T_x\overline G\) have the same dimension.}

\item[\SGAInlineRefTarget{prop:sga3-VIA-1.3.1:item:2}(2)] \textit{The following assertions are equivalent:}
  \begin{enumerate}
  \item[(i)] \textit{\(G\otimes_A K\) is reduced.}
  \item[(i bis)] \textit{\(\mathcal O_{G,e}\otimes_A K\) is reduced.}
  \item[\SGAInlineRefTarget{prop:sga3-VIA-1.3.1:item:ii}(ii)] \textit{\(G\) is smooth over \(A\).}
  \item[(ii bis)] \textit{\(G\) is smooth over \(A\) at the origin.}
  \end{enumerate}
\end{enumerate}

\noindent\textit{Proof.}
\SGARefLink{concl:sga3-VIA-1.3.1-1}{(1)} Let \(x\) be a closed point of \(\overline G\). There is exactly one
\(\overline k\)-morphism
\[
  s:\operatorname{Spec}\overline k\longrightarrow\overline G
\]
whose image is \(x\). Right translation \(r_s\) then induces an
isomorphism from
\[
  \mathcal O_{\overline G,e}
    =\mathcal O_{G,e}\otimes_A\overline k
\]
onto \(\mathcal O_{\overline G,x}\), which proves~\SGARefLink{concl:sga3-VIA-1.3.1-1}{(1)}.

We now prove~\SGARefLink{prop:sga3-VIA-1.3.1:item:2}{(2)}. By SGA~1, II.2.1, we reduce at once to the case in which
\(A\) is a field (\(A=k\)). The implications
\[
  \text{(i)}\Rightarrow\text{(i bis)},\qquad
  \text{(ii)}\Rightarrow\text{(ii bis)},\qquad
  \text{(ii)}\Rightarrow\text{(i)},\qquad
  \text{(ii bis)}\Rightarrow\text{(i bis)}
\]
are evident, so it remains only to prove that (i bis) implies~\SGARefLink{prop:sga3-VIA-1.3.1:item:ii}{(ii)}.

Now (i bis) implies that
\(\mathcal O_{G,e}\otimes_k\overline k\) is reduced. Hence, by~\SGARefLink{concl:sga3-VIA-1.3.1-1}{(1)},
\(\mathcal O_{\overline G,x}\) is reduced for every closed point \(x\) of
\(\overline G\), and therefore \(\overline G\) is reduced. Since
\(\overline G\) is locally of finite type over \(\overline k\), there is
at least one closed point \(y\) such that
\(\mathcal O_{\overline G,y}\) is regular. By~\SGARefLink{concl:sga3-VIA-1.3.1-1}{(1)}, the local rings at all
closed points of \(\overline G\) are isomorphic to
\(\mathcal O_{\overline G,e}\); thus all these local rings are regular.
Consequently, \(\overline G\) is smooth over \(\overline k\), and hence
\(G\) is smooth over \(k\).

\footnote{\label{note:sga3-VIA-15}Editor's note: the following sentence
and \hyperref[ex:sga3-VIA-1.3.2]{Examples~1.3.2} have been added.}
We may now give the examples below, pointed out by M.~Raynaud, of group
schemes \(G\) over a nonperfect field \(k\) for which \(G_{\mathrm{red}}\)
is not a \(k\)-group scheme.
